


Lorsque @@M0@@ est propre sur @@M1@@, on sait [4] que @@M2@@ est représentable par un schéma en groupes localement de type fini @@M3@@ sur @@M4@@. La composante neutre @@M5@@ de @@M6@@ est un groupe algébrique commutatif connexe sur @@M7@@.

\textbf{Définition 2.4.1.} \uline{Soit} @@M8@@ \uline{un schéma propre sur un corps} @@M9@@, @@M10@@ \uline{un} @@M11@@\uline{-module inversible} (resp. @@M12@@ \uline{un élément de} @@M13@@\uline{). On dit que} @@M14@@(resp. @@M15@@\uline{) est} @@M16@@\uline{-équivalent à zéro, s'il existe un entier} @@M17@@ \uline{tel que} @@M18@@(resp. @@M19@@\uline{) est algébriquement équivalent à zéro. Les éléments de} @@M20@@ \uline{algébriquement équivalents} (resp. @@M21@@\uline{-équivalents) à zéro forment un sous-groupe} @@M22@@ (resp. @@M23@@\uline{) de} @@M24@@.

\textbf{Lemme 2.4.3.} \uline{Pour tout schéma propre} @@M25@@ \uline{sur un corps algébriquement clos} @@M26@@ \uline{le groupe} @@M27@@ \uline{est divisible par tout entier} @@M28@@ \uline{premier à la caractéristique de} @@M29@@.

\emph{Preuve} : Comme @@M30@@ est algébriquement clos, on a @@M31@@. Puisque @@M32@@ est connexe et de type fini, l'endomorphisme puissance @@M33@@-ème est surjectif lorsque @@M34@@ est premier à la caractéristique de @@M35@@ (SGA 3 XV 1.3 et VI 1.3.2). Mais @@M36@@ est algébriquement clos, donc l'endomorphisme correspondant sur le groupe @@M37@@ est surjectif, C.Q.F.D.

De 2.4.3 on conclut que dans la suite exacte
@@M38@@
@@M39@@ est divisible par un nombre premier @@M40@@ et @@M41@@ est un groupe de torsion, de sorte qu'on peut appliquer 2.1.3 à la restriction de l'application polynôme (2.3.1) à @@M42@@. On trouve ainsi :

